\section{Design}
\label{sec:design}
A rope is a binary tree whose leaves hold short strings. Reading the leaves
from left to right yields the whole document. Each node $v$ carries a
weight $w(v)$, the number of characters in its left subtree:
\begin{equation}
  w(v) = |\mathrm{left}(v)|.
  \label{eq:weight}
\end{equation}
Weights are the only extra information a rope needs. Every operation below
reads them on the way down the tree and updates them on the way back up,
so keeping them correct costs nothing beyond the nodes an operation already
touches.

\subsection{Finding a position}
\label{sec:finding}
To find the character at position $i$, we start at the root and compare $i$
with the weight of the current node. If $i < w(v)$, we continue in the left
subtree; otherwise we continue in the right subtree with $i - w(v)$. When we
reach a leaf, $i$ is an offset into its string. A rope of depth $d$ answers
this query in $O(d)$ steps.

We call a rope balanced when its depth $d$ and its length $n$ satisfy
\begin{equation}
  n \ge F_{d+2},
  \label{eq:balance}
\end{equation}
where $F_k$ is the $k$-th Fibonacci number~\cite{boehm}. Since $F_k$ grows
like $\varphi^k$ with $\varphi = (1+\sqrt{5})/2$, condition~\eqref{eq:balance}
bounds the depth of a balanced rope by $\log_\varphi n + O(1)$, so finding a
position takes logarithmic time.

\subsection{Editing}
\label{sec:editing}
Splitting walks down the same path as finding a position and creates new
nodes only along that path. Concatenation creates a single new root, whose
weight is the length of its left argument.
In order to be able to insert a string at position $i$, we first split the rope at $i$ into two ropes, and then we concatenate the left part, the new string, and the right part.
Neither operation copies the text stored in the leaves, which is why ropes
stay fast for large documents.

\begin{figure}[t]
  \centering
  \fbox{split at $i$} \enspace$\rightarrow$\enspace
  \fbox{left part} \enspace$+$\enspace \fbox{new text} \enspace$+$\enspace \fbox{right part}
  \enspace$\rightarrow$\enspace \fbox{concatenate}
  \caption{Inserting text into a rope.}
  \label{fig:insertion}
\end{figure}


Each edit therefore costs a split and at most two concatenations. Over a
sequence of $k$ edits, where the document has $n_j$ characters before the
$j$-th edit, the total cost is
\begin{align}
  T(k) &= \sum_{j=1}^{k} O(\log n_j) \label{eq:cost} \\
       &\le k \cdot O(\log n_k).
\end{align}
The bound assumes that the rope stays balanced between edits.

Deletion works in the same way: split at both ends of the deleted range and
concatenate the outer parts. Repeated edits can unbalance the tree, so an
editor occasionally rebuilds it, which restores condition~\eqref{eq:balance}
in time linear in the number of leaves.
Rebalancing is cheap enough to run after every few hundred edits.

The size of the leaves is a tuning parameter. Very short leaves waste memory
on internal nodes, while very long leaves make every split copy more text
than necessary. Leaves of a few hundred bytes work well in practice, because
they fit in a handful of cache lines and keep the tree shallow even for
documents of many megabytes.

Because old nodes are never modified, a rope also gives the editor an undo
history almost for free: each version of the document is simply a pointer
to an older root.
